\section{God heals the heart that he commands to love}

A whole heart is more than a heart stirred by a strong religious feeling. In the Gospel it names the undivided direction of a person's love, reaching thought, desire, and action. Augustine explains the command as leaving no part of the self outside its relation to God. Even love of oneself finds its proper order there. A person who loves himself as though he were his own final good misunderstands what is good for him; love becomes rightly ordered when it seeks the God in whom its good is found. The neighbor belongs within that ordering, because the good sought for oneself is also to be sought for another.\footnote{Augustine, \textit{De doctrina christiana}, I.22.20--21.}

That account exposes a difficulty. A hearer may recognize the order of love, approve it, and still discover reluctance when an actual neighbor makes a claim. The Gospel's scope exceeds the ease of liking pleasant people or enjoying worship when untroubled. The appointed prayers answer this difficulty without making the command smaller. The Collect asks God to free the people from corruption and give a pure following; the Secret asks the holy action to reach past offenses and future dangers; the Postcommunion asks for the cure of vices. What is demanded in the Gospel is sought as healing throughout the Mass.

\subsection{Mercy teaches more than the correct answer}
Augustine's reading of the Introit's mercy verse identifies the decisive dependence. The psalmist does not ask to be treated according to his own righteousness. He asks God to teach the righteousness by which God makes persons righteous. The same psalm is full of promised obedience, so its petitions cannot mean that obedience is unnecessary. Rather, the willing servant continually receives the light and strength by which he serves. Augustine's next paragraph insists that understanding must continue to be received: a former illumination does not render the present person independent of the source of light.\footnote{Augustine, \textit{Psalms}, Ps.~118, English paragraphs 123--124; compare 34--36.}

Bellarmine draws out the difference between knowledge alone and practical love. Teaching at verse 124 includes the conviction and affection that make observance desirable. This is especially searching beside the lawyer's question. One can know that love is the greatest command and use that knowledge to gain superiority over someone else. Then the correct answer has not yet become the life it describes. God's teaching reaches further than the acquisition of a sentence: it makes that sentence a rule for the speaker's own conduct.\footnote{Bellarmine, \textit{Psalms}, Ps.~118, commentary on v.~124.}

The Introit's first words prevent mercy from becoming evasion. God is just, and his judgment is right. To ask for mercy is therefore to relinquish the attempt to make one's wrong right by argument. The servant's need can be admitted without despair because the judge is addressed as the giver of mercy. The accompanying verse's blessed walkers show the desired outcome. The worshipper is not praying to remain indefinitely at the point of failure; he asks to join a way of life that can be walked.

The Collect carries that movement into allegiance. A pure mind follows God alone because competing masters corrupt its direction. This does not divide human life into sacred occasions that belong to God and ordinary occasions in which another rule applies. The Gospel names the whole self; Paul's virtues name its ordinary conduct. A desire to appear devout can coexist with impatient treatment of another person, but the texts do not allow that coexistence to remain comfortable.

\subsection{The neighbor tests the heart's unity}
Chrysostom connects the love commands by their consequences. To love God is to attend to what God commands; the command includes the neighbor. Conversely, hatred harms the very love of God a person may profess. The neighbor is therefore not an optional second project after a sufficient amount of religious attention has been given to God. The second command reveals whether the first is being lived.\footnote{Chrysostom, \textit{Matthew}, hom.~71, exposition of Matt.~22:37--40.}

Augustine's treatment of neighbor intensifies the test by including the enemy. An individual's worthiness of affection does not determine whether he is human or whether mercy may be owed to him. Augustine also brings Christ within the name of neighbor through the mercy Christ has shown us. The two directions belong together: Christ gives mercy, and a recipient of that mercy is called to act mercifully. Loving the neighbor in God does not turn the neighbor into a tool. It seeks his true good, including the good that selfish indulgence could deny him.\footnote{Augustine, \textit{De doctrina christiana}, I.30.31--33.}

The Epistle shows how that love becomes visible when goodwill is tested. Humility gives up the need to be superior; meekness restrains the impulse to hurt; patience bears a burden over time; charity keeps endurance directed toward the other's good. None requires delight in another person's faults. Thomas says that bearing with someone includes correction at a suitable time. The difference between patient charity and negligent agreement is the good one seeks for the other. Correction loses that purpose when it becomes revenge; patience loses it when it becomes indifference.\footnote{Thomas, \textit{Super Ephesios}, ch.~4, lect.~1, Dessain vol.~2, pp.~312--315.}

This makes healing observable without reducing it to a performance. A concrete fruit might be refusing the humiliating remark one is capable of making, telling a needed truth without contempt, or continuing a responsibility after enthusiasm has waned. These are applications of Paul's virtues, not promises of immediate emotional ease. A person can act in charity while still needing healing in desire. The prayer and the practice belong to the same conversion.

\subsection{The Lord who gives what he asks}
The second half of the Gospel prevents the first from being reduced to advice from an impressive moral teacher. Chrysostom reads Jesus' question as a disclosure of divinity: Davidic descent is affirmed, while the reduction of Christ to a merely human descendant is overturned by David's own confession. The Lord is not absent while the hearer tries to fulfill his command. The one to whom love is owed enters the question of how love is given.\footnote{Chrysostom, \textit{Matthew}, hom.~71, exposition of Matt.~22:41--46.}

The Gradual supplies a fitting image of dependence. The heavens stand by the Word, and their strength comes through the Spirit. Augustine identifies spiritual heavens in the righteous, established by that same Word rather than by themselves. The image does not make human obedience unreal. Firmness is real precisely as a received firmness. The blessed people of the first verse likewise exists by election before it can offer itself as an achievement.\footnote{Augustine, \textit{Psalms}, Ps.~32, first exposition, paragraphs 6 and 12.}

The Alleluia allows a needy heart to speak instead of waiting until its need has vanished. Bellarmine's account of grace-enabled prayer means that the first honest cry already turns toward the physician. At the Offertory that cry expands to the sanctuary and the people marked by God's name. Daniel's plea resists a narrow project of private improvement. The heart being healed belongs to a people whose wounds it must learn to carry in intercession. The clean mind of the Collect is thus compatible with grief over a desolate sanctuary: purity need not become distance from the afflicted.\footnote{Bellarmine, \textit{Psalms}, Ps.~101, opening commentary; Dan.~9:17--19; Jerome, \textit{Daniel}, II, on 9:17--18.}

Thomas locates the efficacy of the Eucharist in Christ and his Passion. Spiritual nourishment is not the power of an encouraging idea alone. The sacrament sustains and increases the life of grace because Christ is its source. That teaching illuminates the Secret's request concerning the holy action and the Postcommunion's language of cure. The mysteries do what the wounded person cannot achieve by merely condemning his own wound.\footnote{Thomas, \textit{Summa}, III, q.~79, a.~1, corpus.}

Yet preservation does not abolish freedom. Thomas explicitly allows the possibility of later sin. The Secret cannot be read as immunity from responsibility, and reception cannot be made a guarantee that no serious refusal will follow. The Communion's demand for vows to be paid gives the healed person a task, while its feared Lord denies him mastery over judgment. Augustine's account joins the command to God's enabling help: fidelity is required, and reliance on one's unaided strength will fail.\footnote{Thomas, \textit{Summa}, III, q.~79, a.~6, corpus and ad 1; Augustine, \textit{Psalms}, Ps.~75, paragraph 11.}

\subsection{The four senses of healed charity}
\textbf{Literal.} The formulary joins distinct acts of speech: the psalmist's obedience and supplication, Paul's appeal to a called community, Jesus' commands and question, Daniel's intercession, and the Church's petitions for purification. The law requires love; the prayers ask for real help. Neither command nor need disappears when they are heard together.

\textbf{Allegorical.} David's Son is David's Lord, and the divine Word strengthens his people through the Spirit. Christ meets the afflicted in his body and gives sacramental life through his Passion. The healing sought by the prayers reaches its Christian center in the Lord whose command the Gospel proclaims.

\textbf{Moral.} The whole person learns to direct love toward God and to seek the neighbor's good, including an enemy's. Humility, patient endurance, fitting correction, and fulfilled commitments give that love its daily form. Asking for grace and acting with grace are one movement, rather than competing explanations of obedience.

\textbf{Anagogical.} The cure of vice tends toward the eternal remedy. Thomas's account of perfect peace and unity in glory gives the Postcommunion its full horizon: love's healing is completed in life with God and the saints. Judgment remains serious because the heart's direction matters; hope remains possible because the remedy is God's gift.\footnote{Thomas, \textit{Summa}, III, q.~79, a.~2, corpus.}
